\begin{figure}[!ht]
	\tikzsetnextfilename{fig_ctrl_torque_inner_loop}
	\centering
	
	\tikzstyle{comp} = [circle, draw, minimum size = 0.5cm, line width=0.7] % comparator style
	\tikzstyle{block} = [rectangle, draw, minimum width = 1.2cm, minimum height = 1.2cm, line width=0.7]  % bloc style
	
	\newcommand{\comp}[3][]{ % comparator command
		\node[comp] (#2) at #3 {};
		\draw (#2) -- + (45:0.25) -- + (45:-0.25);
		\draw (#2) -- + (-45:0.25) -- + (-45:-0.25);
		\ifthenelse{
			\isempty{#1}{}
		}{}
		{
			\ifthenelse{
				\equal{#1}{a}
			}{
				\node[above right] at (#2) {$-$\phantom{p}};
			}{}
			\ifthenelse{
				\equal{#1}{b}
			}{
				\node[below right] at (#2) {$-$\phantom{l}};
			}{}
		}
	}
	\resizebox{\textwidth}{!}{
		\begin{tikzpicture} [node distance = 2cm]
			
			\comp[b]{ci}{(0, 0)};
			\node[block, right of = ci, align = center] (current) {PID\\courant};
			\node[block, right of = current] (pwm) {MLI\footnotemark};
			\node[block, right of = pwm] (motor) {Moteur};
			\node[block, left of = ci, align = center] (velocity) {PID\\vitesse};
			\node[block, left of = velocity, align = center, node distance = 2.5cm] (rampGen) {Générateur\\de rampe};
			\comp[b]{cv}{($ (rampGen) + (-2,0) $)};
			%
			\node[block, below of = motor, align = center, node distance = 1.4cm] (currentSensor) {Capteur de\\courant};
			\node[block, below of = currentSensor, align = center, node distance = 1.4cm] (velocitySensor) {Capteur de\\vitesse};			
			
			\draw[-{Latex[length=3mm]}] (cv.east) -- (rampGen.west);
			\draw[-{Latex[length=3mm]}] (rampGen.east) -- (velocity.west);
			\draw[-{Latex[length=3mm]}] (velocity.east) -- (ci.west)
				node[midway, above]{$i_d$};
			\draw[-{Latex[length=3mm]}] (ci.east) -- (current.west);
			\draw[-{Latex[length=3mm]}] (current.east) -- (pwm.west);
			\draw[-{Latex[length=3mm]}] (pwm.east) -- (motor.west);
			
			\draw[-{Latex[length=3mm]}] (currentSensor.west) -| (ci.south);
			\draw[-{Latex[length=3mm]}] (velocitySensor.west) -| (cv.south);
			\draw[-{Latex[length=3mm]}] ($(cv.west) + (-1.25,0)$) -- (cv.west)
				node[midway, above]{$\dot{q}_d$};
			
		\end{tikzpicture}
	}
	\caption{Schéma de contrôle cascade proposé par les pilotes moteurs implantés sur le youBot, pour un moteur \parencite{TRINAMIC_manual_2011}, et une régulation en vitesse}
	\label{fig_trinamic_cascade_ctrl_loop}
\end{figure}
\footnotetext{MLI: modulation à largeur d'amplitude, aussi connue sous l'acronyme anglais PWM pour \textit{pulse width modulation}.}